% docs/style/uni-makros.tex -- local copy of tex-skripten/formatierungen/uni-makros.tex
% Copied 2026-10-12 so that the repository compiles stand-alone (GitHub).
% Local version: 1.0.0 -- unchanged copy



% bis zur Stufe 6 werden Überschriften nummeriert
\setcounter{tocdepth}{6}
\setcounter{secnumdepth}{6} % Kapitelnummerierung bis ...

% Zeilenabstand innerhalb Tabellen vergroessern
\renewcommand{\arraystretch}{1,5}

% Schusterjungen und Hurenkinder vermeiden (10000 = unendlich)
\clubpenalty = 10000
\widowpenalty = 10000
\displaywidowpenalty = 10000

%====================================================
% Kopf- und Fußzeilen einrichten
%====================================================
%\clearpairofpagestyles			% default ausschalten


%---------------------------------------------------


% Schrift Helvetica im gesamten Dokument:
\def\rmdefault{phv}
\def\sfdefault{phv}
\def\ttdefault{pcr}

% damit Listings syntaxgefärbt ausgegeben werden
\definecolor{Brown}{cmyk}{0,0.81,1,0.60}
\definecolor{OliveGreen}{cmyk}{0.64,0,0.95,0.40}
\definecolor{CadetBlue}{cmyk}{0.62,0.57,0.23,0}
\definecolor{Orange}{cmyk}{0.5,0.7,0.9,0}
\definecolor{lightblue}{rgb}{0.8,0.85,1}
\definecolor{rahmen}{rgb}{0.9,0.9,0.7}
\definecolor{syntax}{cmyk}{0.0,0.0,0.8,0.0}

\lstset{frame=ltrb,framesep=5pt,basicstyle=\small,breaklines=true,
	keywordstyle=\ttfamily\color{OliveGreen},numbers=left,
	identifierstyle=\ttfamily\color{CadetBlue}\bfseries,commentstyle=\color{Brown},
	stringstyle=\ttfamily\color{Orange}}

%====================================================
% für nummerierte Beispiele
% Klammern mit \begin{bsp}  \end{bsp}
%====================================================
\newtheorem{defi}{Definition}[chapter]
\newtheorem{bsp}[defi]{Beispiel}

%====================================================
% Farbpalette für TikZ-Diagramme (Skriptum Teil I-III)
% Wird in allen tex/teil*.tex-Dateien für Diagramme verwendet
%====================================================
\definecolor{myblue}{HTML}{185FA5}
\definecolor{myteal}{HTML}{0F6E56}
\definecolor{mycoral}{HTML}{993C1D}
\definecolor{myamber}{HTML}{854F0B}
\definecolor{mypink}{HTML}{D4537E}

%====================================================
% Makro für Geschwindigkeits-/Ensemblemittel (Teil III.b, III.c)
% \avg{v_i} entspricht \overline{v_i}
%====================================================
\newcommand{\avg}[1]{\overline{#1}}

%====================================================
% Makros zur Vereinfachung
%====================================================
\newcommand{\Index}[1]{#1\index{#1}}

\def\TODO#1{\color{blue}{\textbf{TODO: #1}}\color{black}}

\def\id{
	\begin{flushright}
		\vspace{-10mm}
		\begin{tiny}ID: \currfilename\end{tiny}
	\end{flushright}
}


\newcommand\prozent[1]{\hfill\textbf{#1}\%}
%\newcommand\Laplace[1]{\begin{Large}$\mathcal{L}$\end{Large}{#1}}

%Grad-Zeichen
%\newcommand\grad{$^\circ$}

%entspricht-Zeichen.
\newcommand{\entspricht}{\ensuremath{\mathrel{\widehat{=}}}}

\newcommand\circlearound[1]{%
	\tikz[baseline]\node[draw,shape=circle,anchor=base] {#1} ;}
%====================================================
% Bilder einfügen (zitat mit biber)
%====================================================
\def\abb#1#2#3#4#5{% 				Grösse,Dateiname,Beschriftung,Label,Bildzitatquelle
	\begin{figure}[ht!]
		\center
		\includegraphics[#1]{#2}
		\begin{flushright}
			\vspace{-7mm}
			\begin{tiny}ID: #2\end{tiny}
		\end{flushright}
		\vspace{-7mm}
		\caption{#3}(Quelle:\autocite{#5})
		\label{#4}
	\end{figure}
}




\def\tikzbild#1#2#3#4{% 				scaling-factor,Dateiname,Beschriftung,Label
	\begin{figure}[h!]%
		\centering
		\scalebox{#1}{
			\input{#2}%
		} %end of scalebox
		\begin{flushright}%
			\vspace{-7mm}%
			\begin{tiny}ID: #2\end{tiny}%
		\end{flushright}%
		\vspace{-7mm}%
		\caption{#3}%
		\label{#4}(Quelle:\cite{pic:scub-tikz})%
	\end{figure}
}



%====================================================
% Tabelle einfügen
%====================================================
\def\tabelle#1#2#3#4#5{% 			Spaltenaufteilung,Beschriftung,Label,fontsize,Tabelle
	\begin{table}[!ht]
		\begin{center}#4
			\begin{tabular}{#1}
				#5
			\end{tabular}
		\end{center}
		\caption{#2}
		\label{#3}
	\end{table}
}

%------------------ Farbbox -------------------------
% Verwendung:
%\begin{rahmenbox}[colback=green!5,colframe=blue!50]{Titel}
%Box-Text
%\end{rahmenbox}


\newenvironment{rahmenbox}[2][]{%
	\tcolorbox[savedelimiter=rahmenbox, savelowerto=\jobname_bspsave2.tex, lowerbox=ignored,fonttitle=\bfseries,title=#2, #1]}%
{\endtcolorbox}


%====================================================
% Datenblatt (pdf) einfügen
%====================================================
\def\datenblatt#1#2{% 			Dateiname ohne .pdf, Seiten von-bis
	\index{Datenblatt #1}
	\includepdf[pages=#2,pagecommand={},scale=0.7]{#1.pdf}
}

%====================================================
% gelbe equation.box \begin{tcolorbox}  \end{tcolorbox}
%====================================================
%\tcbset{colback=yellow!10!white, colframe=red!50!black,
	%	highlight math style= {enhanced, %<-- needed for the ’remember’ options
		%		colframe=red,colback=red!10!white,boxsep=0pt}
	
	%====================================================
	% Astrozeichen
	%====================================================
	%\newcommand\circlearound[1]{%
		%	\tikz[baseline]\node[draw,shape=circle,anchor=base] {#1} ;}
	%====================================================
	\newcommand*{\Sun}{{\odot}}
	\newcommand*{\Mercury}{{\text{\mathabxbfamily\char"41}}}
	\newcommand*{\Venus}{{\text{\mathabxbfamily\char"42}}}
	\newcommand*{\Earth}{{\text{\mathabxbfamily\char"43}}}
	\newcommand*{\Mars}{{\text{\mathabxbfamily\char"44}}}
	\newcommand*{\Jupiter}{{\text{\mathabxbfamily\char"45}}}
	\newcommand*{\Saturn}{{\text{\mathabxbfamily\char"46}}}
	\newcommand*{\Uranus}{{\text{\mathabxbfamily\char"47}}}
	\newcommand*{\Neptune}{{\text{\mathabxbfamily\char"48}}}
	\newcommand*{\Pluto}{{\text{\mathabxbfamily\char"49}}}
	
	%Den Punkt am Ende jeder Beschreibung deaktivieren
	\renewcommand*{\glspostdescription}{}
	
	%\newcommand{\mygls}[1]{\glsadd{#1} \gls{#1}...\glsentrydesc{#1}}
	%\newcommand{\mygls}[1]{\gls{#1}:~\glsentrydesc{#1}}
	%\newcommand{\mygls}[1]{\glsadd{#1}\glsentryname{#1}:~\glsentrydesc{#1}}
	\newcommand{\abk}[1]{%
		\ifglsused{#1}{}{%
			\glsadd{#1}%
		}%
		\gls{#1}:~\glsentrydesc{#1}%
	}
	\newcommand{\mygls}[1]{%
		\ifglsused{#1}{}{%
			\glsadd{#1}%
		}%
		\gls{#1}:~\glsentrydesc{#1}%
	}
	
	
	% sinnvoll, aber sort in den def muesste angepasst werden.....
	\newcommand{\glssection}[2]{%
		\ifglsused{#1}{}{%
			\newglossaryentry{#1}{
				type=formel,
				name={{\Large \textbf{#2}}},
				symbol={},
				description={},
				sort=#1-0000
			}%
			\glsadd{#1}%
		}%
	}
	
	
	\newcommand{\tabcite}[1]{%
		\cite{#1} & \citeauthor{#1} & \citetitle{#1}%
	}
	%====================================================
	% Einheiten für siunitx
	%====================================================
	\DeclareSIUnit\parsec{pc}
	\DeclareSIUnit\lightyear{ly}
	
	
	
	
	% --- Definition für Code-Blöcke (Bash/Shell) ---
	\lstdefinestyle{bashstyle}{
		language=bash,
		backgroundcolor=\color{black!5},
		commentstyle=\color{green!40!black},
		keywordstyle=\color{blue},
		numberstyle=\tiny\color{black!50},
		stringstyle=\color{purple},
		basicstyle=\footnotesize\ttfamily,
		breaklines=true,
		captionpos=b,
		keepspaces=true,
		numbers=none,
		numbersep=5pt,
		showspaces=false,
		showstringspaces=false,
		showtabs=false,
		tabsize=2,
		frame=single,
		rulecolor=\color{black!20},
		keywords={python, python3, pip, source, deactivate, m, venv, install, freeze, r}
	}
	
	% --- Definition für PowerShell ---
	% KORRIGIERT: 'language=PowerShell' zu 'language=bash' geändert.
	% 'PowerShell' ist keine Standard-Sprache im 'listings'-Paket.
	% 'bash' dient als robuster Ersatz für die Syntax-Hervorhebung.
	\lstdefinestyle{psstyle}{
		language=bash,
		backgroundcolor=\color{black!5},
		commentstyle=\color{green!40!black},
		keywordstyle=\color{blue},
		basicstyle=\footnotesize\ttfamily,
		breaklines=true,
		captionpos=b,
		numbers=none,
		frame=single,
		rulecolor=\color{black!20},
		keywords={python, pip, Set-ExecutionPolicy, RemoteSigned, Scope, CurrentUser, install, freeze, r, .\.venv\Scripts\Activate.ps1}
	}
